\uuid{AIt8}
\titre{Continuité en un point}
\auteur{bodin ; hanani ; mzari}
\organisation{exo7}
\licence{CC BY-NC-SA 4.0}
\datecreate{2020-01-20}
\niveau{L1}
\module{Analyse}
\chapitre{Continuité, limite et étude de fonctions réelles}
\sousChapitre{Continuité : théorie}
\difficulte{2}
\theme{Fonctions continues, continuité en un point}
\exo7qcmid{lille-1.continuite.014}

\contenu{
\begin{question}
\qclassification{123.01,123.02}
\qtags{motcle=continuité en un point}
Parmi les propriétés suivantes, quelles sont celles qui implique que $f$ est continue en $x_ 0$ ?
\begin{answers}
    \good{$\lim_{x\to x_0} f(x) = f(x_0)$}

    \good{$\forall \epsilon >0 \quad \exists \delta > 0 \qquad
    |x-x_0| \le \delta \implies |f(x)-f(x_0)| \le \epsilon$}

    \bad{$\exists \delta > 0 \quad \forall \epsilon >0 \qquad
    |x-x_0| < \delta \implies |f(x)-f(x_0)| < \epsilon$}

    \good{$\big| f(x) - f(x_0) \big| \to 0$ lorsque $x \to x_0$}

\end{answers}
\begin{explanations}
$f$ est continue en $x_0$ si $\lim_{x\to x_0} f(x) = f(x_0)$, ce qui s'écrit aussi $\big| f(x) - f(x_0) \big| \to 0$, ou encore :
$\forall \epsilon >0 \quad \exists \delta > 0 \qquad
    |x-x_0| < \delta \implies |f(x)-f(x_0)| < \epsilon$ (et on peut remplacer les inégalités strictes par des inégalités larges).
\end{explanations}
\end{question}
}
